\subsection{Assumptions}
Table~\ref{tab:assumptions} states which assumptions each result requires. We separate the characterization of exact encoders from the stronger conditions needed for a positive sharp separation score. Differentiability only in the interior of $I$ does not ensure finite sensitivity at its endpoints. Strict convexity or concavity of the local preprocessing function $\phi$ suffices for the endpoint rank argument used in the score formula; we do not claim that it is necessary.

Throughout, $C^1(I)$ denotes real-valued continuously differentiable functions with finite continuous one-sided derivatives at the endpoints. We state separately when only interior differentiability is needed.

\begin{table*}[t]
\centering
\setlength{\tabcolsep}{5pt}
\caption{Assumptions and scope of the main results. Throughout, $I=[-1,1]\subset\R$, $K\ge4$, $P>0$, $f:I^K\to\R$, and the shared encoder $\enc:I\to\R^2\cong\C$ is continuous. Optimization results impose the pointwise power constraint \eqref{eq:peak-model}.}
\label{tab:assumptions}
\footnotesize
\begin{minipage}{0.98\textwidth}
\textbf{R (exact-encoder characterization):} $\phi,\psi\in C(I)\cap C^1((-1,1))$; $\phi$ is nonaffine; $\psi$ is strictly convex or strictly concave.\\[2pt]
\textbf{S (sharp separation and design):} R, with $\phi$ strictly convex or strictly concave and $\psi\in C^1(I)$.
\end{minipage}
\vspace{5pt}
\renewcommand{\arraystretch}{1.2}
\begin{tabular}{@{}>{\raggedright\arraybackslash}p{0.27\textwidth}>{\centering\arraybackslash}p{0.07\textwidth}>{\raggedright\arraybackslash}p{0.60\textwidth}@{}}
\toprule
Result & Regime & Scope and additional requirements\\
\midrule
Exact encoders, Theorem~\ref{thm:rigidity} & R & $f=f_{\phi,\psi}$; exactness as in \eqref{eq:exactness}.\\
Optimal separation, Theorem~\ref{thm:sdp} & S & $f=f_{\phi,\psi}$; global optimization over $\mathcal E_P$.\\
Explicit curves, Theorem~\ref{thm:parabola} & S & $\phi(s)=s^2$ or $s^4$; $\psi$ is even.\\
Two-statistic extension, Theorem~\ref{thm:extension} & --- & $\gamma\in C(I;\R^2)\cap C^1((-1,1);\R^2)$; $1,\varphi_1,\varphi_2$ linearly independent; $F\in C(\mathcal M;\R)$; condition~\eqref{eq:completion-id}. Part 2 also requires nonconstant Lipschitz $F$. Formula~\eqref{eq:polyopt} further assumes quadratic $\gamma$, $C^1$ postprocessing near $\mathcal C_2$, and \eqref{eq:general-endpoints}.\\
Projection receiver, Theorem~\ref{thm:decoder} & S & $f=f_{\phi,\psi}$; $\enc(s)=c+us+v\phi(s)$ with $\det[\,u\ v\,]\ne0$; specified strip-projection receiver.\\
Small-noise frontier, Theorem~\ref{thm:pareto} & S & Fixed joint input distribution $\pi$; specified strip decoder; $|m|<1$ almost surely; $Q_\pi\succ0$; independent circular Gaussian noise; any output clipping is locally inactive almost surely.\\
Certified design, Theorem~\ref{thm:certifiedgrid} & S & $\phi\in C^2(I)$; grids contain both endpoints and three fixed points with linearly independent $z(s)$; numerical implementation uses verified function, derivative, and matrix bounds.\\
\bottomrule
\end{tabular}
\end{table*}

\paragraph{Endpoint slopes}
Consider $\phi(s)=s^2$ and $\psi(s)=-\sqrt{1-s}$. Both satisfy R, but $\psi'$ diverges at $1$. Theorem~\ref{thm:rigidity} makes every exact continuous encoder an affine quadratic curve. Compare the two equal-input configurations $s_k=1$ and $s_k=1-h$ for all $k$, where $h\downarrow0$. Their received difference is
$$
\Delta r=K[-hu+(-2h+h^2)v]=O(h),
$$
whereas the difference between their desired function values is $\sqrt h-2h+h^2$. Their squared-separation quotient therefore tends to zero. Every exact encoder has zero score, and every nonexact encoder has zero score by \eqref{eq:score}. Thus R alone permits a zero global optimum and no receiver with finite uniform disturbance amplification.

We therefore first establish which sums exactness requires the transmitters to preserve, and then impose S to optimize robustness.
